% Relativistic Chapter VII, §64–66: Rayleigh Criterion for Couette Flow
% Relativistic generalization of Chandrasekhar, Chapter VII §64–66
% Conventions: see RELATIVISTIC_CONVENTIONS.md

\chapter{The Stability of Relativistic Couette Flow}
\label{ch:rel-7}

%% ================================================================
\section{Introduction: relativistic rotating flows}\label{sec:rel-7-64}

In Chapter~VII of the non-relativistic treatment, the stability of
Couette flow---the steady circular flow of a liquid between two
rotating coaxial cylinders---was analysed and found to hinge on the
radial gradient of specific angular momentum.  The present chapter
extends that analysis to the relativistic regime, where the inertia
of energy plays a decisive role.

Relativistic rotation arises naturally in at least three
astrophysical contexts:
\begin{enumerate}
\item \textbf{Accretion disks.}  Matter spiralling into a compact
  object (neutron star or black hole) can reach orbital velocities
  $v_{\phi}\sim 0.1$--$0.5\,c$.  The inner disk is effectively a
  relativistic Couette flow, and its stability governs angular-momentum
  transport and, ultimately, the observed luminosity.
\item \textbf{Relativistic jets.}  The boundary layer between a
  fast-moving jet ($\Gamma\gg 1$) and the ambient medium experiences
  strong rotational shear.  Whether that shear layer is Rayleigh-stable
  determines the growth of Kelvin--Helmholtz modes and the jet's
  long-range collimation.
\item \textbf{Neutron-star interiors and heavy-ion collisions.}
  The quark--gluon plasma produced in collisions of heavy nuclei, and
  the dense matter in differentially rotating neutron stars, can sustain
  rotational velocities for which $(\varepsilon+p)/c^{2}$ differs
  appreciably from the rest-mass density $\rho_{0}$.
\end{enumerate}

The non-relativistic criterion of Rayleigh states that a
distribution of angular velocity $\Omega(r)$ is stable if and only if
the square of the specific angular momentum $L=r^{2}\Omega$ increases
outward.  When the flow is relativistic, the ``inertia'' of a fluid
element is no longer its rest-mass density $\rho_{0}$ but the
relativistic enthalpy density $w=(\varepsilon+p)/c^{2}$; the specific
angular momentum must accordingly be redefined.  The purpose of the
next two sections is to make this precise and to derive the
relativistic Rayleigh criterion together with the associated
epicyclic frequency used in modern accretion-disk theory.

Throughout, we work in flat Minkowski spacetime with cylindrical
coordinates $(t,r,\theta,z)$ and signature $(-,+,+,+)$.  The speed
of light $c$ is kept explicit so that the non-relativistic limit is
transparent.  The four-velocity satisfies $u^{\mu}u_{\mu}=-c^{2}$.

%% ================================================================
\section{The physical problem: relativistic Couette flow}
\label{sec:rel-7-65}

\subsection{Steady circular flow}
\label{sec:rel-7-65a}

Consider an ideal (inviscid, non-heat-conducting) relativistic fluid
confined between two coaxial cylinders of radii $R_{1}$ and
$R_{2}$ ($R_{2}>R_{1}$) rotating with angular velocities $\Omega_{1}$
and $\Omega_{2}$.  We seek stationary, axisymmetric, purely circular
flows:
\begin{equation}
\label{eq:rel-7-1}
u^{\mu} = \Lf\bigl(c,\,0,\,r\Omega,\,0\bigr),
\end{equation}
where $\Omega=\Omega(r)$ is the angular velocity and
\begin{equation}
\label{eq:rel-7-2}
\Lf = \frac{1}{\sqrt{1-v_{\phi}^{2}/c^{2}}},\qquad
v_{\phi}=r\Omega,
\end{equation}
is the Lorentz factor.  Causality demands
\begin{equation}
\label{eq:rel-7-3}
\boxed{v_{\phi}=r\Omega < c}
\end{equation}
everywhere in the flow domain.  This is an absolute upper bound that
has no analogue in the non-relativistic theory and restricts the
permissible profiles $\Omega(r)$.

\subsection{Relativistic centrifugal balance}
\label{sec:rel-7-65b}

For a perfect fluid the energy-momentum tensor is
\begin{equation}
\label{eq:rel-7-4}
T^{\mu\nu} = \frac{\varepsilon+p}{c^{2}}\,u^{\mu}u^{\nu}
             + p\,g^{\mu\nu},
\end{equation}
and the equations of motion are $\nabla_{\mu}T^{\mu\nu}=0$.
Projecting perpendicular to $u^{\nu}$ and evaluating for the
steady circular flow \eqref{eq:rel-7-1}, the radial component yields
the \emph{relativistic centrifugal-balance equation}:
\begin{equation}
\label{eq:rel-7-5}
\boxed{\frac{\varepsilon+p}{c^{2}}\;\frac{v_{\phi}^{2}}{r}
       = -\frac{\partial p}{\partial r}\,.}
\end{equation}
Equivalently, introducing the relativistic enthalpy density
$w=(\varepsilon+p)/c^{2}$,
\begin{equation}
\label{eq:rel-7-6}
w\,\frac{v_{\phi}^{2}}{r} = -\frac{\partial p}{\partial r}\,.
\end{equation}
In the non-relativistic limit $\varepsilon\to\rho_{0}c^{2}$,
$p/\varepsilon\to 0$, and~\eqref{eq:rel-7-5} reduces to the familiar
$\rho_{0}\,v_{\phi}^{2}/r = -\partial p/\partial r$.

For a barotropic equation of state $p=p(\varepsilon)$ (or, more
generally, for an isentropic flow), \eqref{eq:rel-7-6} can be
integrated.  Defining the relativistic specific enthalpy by
\begin{equation}
\label{eq:rel-7-7}
h = \int_{0}^{p}\frac{dp'}{\varepsilon(p')+p'}\,,
\end{equation}
one obtains
\begin{equation}
\label{eq:rel-7-8}
h + \tfrac{1}{2}\ln\!\left(1-\frac{v_{\phi}^{2}}{c^{2}}\right) = \text{const},
\end{equation}
which generalises the Bernoulli integral for relativistic rotating flow.

\subsection{Relativistic specific angular momentum}
\label{sec:rel-7-65c}

The azimuthal component of the relativistic Euler equation for the
steady flow~\eqref{eq:rel-7-1}, combined with baryon-number
conservation $\nabla_{\mu}(n\,u^{\mu})=0$, implies that the
\emph{relativistic specific angular momentum},
\begin{equation}
\label{eq:rel-7-9}
\boxed{\ell \;=\; \Lf\,r\,v_{\phi}\;\frac{\varepsilon+p}
               {\rho_{0}\,c^{2}}
       \;=\; \Lf^{2}\,r^{2}\Omega\;\frac{w}{\rho_{0}}\,,}
\end{equation}
is conserved following the fluid, i.e.\ $u^{\mu}\nabla_{\mu}\ell=0$.
Here $\rho_{0}$ is the rest-mass density (baryon mass times number
density) and $w=(\varepsilon+p)/c^{2}$ is the enthalpy density.

The decisive difference from the non-relativistic angular momentum
$L=r^{2}\Omega=r\,v_{\phi}$ is twofold:
\begin{enumerate}
\item The Lorentz factor $\Lf$ appears, reflecting time dilation.
\item The ratio $w/\rho_{0}=(\varepsilon+p)/(\rho_{0}c^{2})$
  weights the angular momentum by the \emph{total inertia per baryon},
  including thermal and binding energy.  In the non-relativistic
  limit $w/\rho_{0}\to 1$ and $\Lf\to 1$, so that $\ell\to L$.
\end{enumerate}

\noindent
\textbf{Remark (isentropic fluid).}
For an isentropic perfect fluid it is often convenient to work with the
\emph{canonical} angular momentum per unit enthalpy,
\begin{equation}
\label{eq:rel-7-10}
\widetilde{\ell} = \Lf\,r\,v_{\phi}
                 = \Lf^{2}\,r^{2}\Omega
                 = \frac{v_{\phi}\,r}{1-v_{\phi}^{2}/c^{2}}\,,
\end{equation}
which is conserved along streamlines for a perfect barotropic fluid
and reduces to $r^{2}\Omega$ in the non-relativistic limit.

%% ================================================================
\section{Rayleigh's criterion: relativistic version}
\label{sec:rel-7-66}

\subsection{Interchange argument}
\label{sec:rel-7-66a}

We follow Rayleigh's original thought experiment, adapted to the
relativistic setting.  Two thin coaxial rings of fluid at radii
$r_{1}$ and $r_{2}$ ($r_{2}>r_{1}$), each of unit rest mass, are
interchanged.  Because angular momentum is conserved, after the
interchange the ring originally at $r_{1}$ carries angular
momentum~$\ell_{1}$ but now sits at~$r_{2}$, and vice versa.

The effective centrifugal potential energy per unit rest mass at
radius~$r$ for a ring of angular momentum~$\ell$ is, in the
relativistic framework,
\begin{equation}
\label{eq:rel-7-11}
E_{\mathrm{cent}}(\ell,r) = \frac{c^{2}}{2}\ln\!\left(
  1+\frac{\ell^{2}}{r^{2}c^{2}}\right),
\end{equation}
which for $\ell\ll rc$ reduces to $\ell^{2}/(2r^{2})$, the
non-relativistic expression.  (This follows from the logarithmic
Bernoulli integral~\eqref{eq:rel-7-8} with $v_{\phi}=\ell/(r\Lf)$.)

The change in total centrifugal potential energy upon interchange is
\begin{equation}
\label{eq:rel-7-12}
\delta E = \bigl[E_{\mathrm{cent}}(\ell_{1},r_{2})
           + E_{\mathrm{cent}}(\ell_{2},r_{1})\bigr]
         - \bigl[E_{\mathrm{cent}}(\ell_{1},r_{1})
           + E_{\mathrm{cent}}(\ell_{2},r_{2})\bigr].
\end{equation}
For infinitesimally separated rings ($r_{2}=r_{1}+\delta r$),
expanding to leading order in $\delta r$ gives
\begin{equation}
\label{eq:rel-7-13}
\delta E \propto \frac{d(\widetilde{\ell}^{\,2})}{dr}\,(\delta r)^{2},
\end{equation}
where $\widetilde{\ell}=\Lf\,r\,v_{\phi}$ is the canonical angular
momentum~\eqref{eq:rel-7-10}.

\emph{If $\delta E>0$, the interchange requires an external source of
energy and the flow is stable; if $\delta E<0$, energy is released
and the flow is unstable.}  Hence:

\subsection{The relativistic Rayleigh criterion}
\label{sec:rel-7-66b}

\begin{quote}
\textbf{Theorem (Relativistic Rayleigh criterion).}
\emph{In the absence of viscosity and heat conduction, a stationary,
axisymmetric, circular flow of a relativistic perfect fluid is stable
to axisymmetric perturbations if and only if}
\begin{equation}
\label{eq:rel-7-14}
\boxed{\frac{d}{dr}\!\left(\widetilde{\ell}^{\,2}\right) > 0
\qquad\text{everywhere},}
\end{equation}
\emph{where $\widetilde{\ell}=\Lf^{2}\,r^{2}\Omega$ is the
relativistic specific angular momentum.  The flow is unstable if
$\widetilde{\ell}^{\,2}$ decreases anywhere in the interval.}
\end{quote}

\noindent
Writing \eqref{eq:rel-7-14} explicitly:
\begin{equation}
\label{eq:rel-7-15}
\frac{d}{dr}\!\left(\frac{r^{4}\Omega^{2}}
   {\bigl(1-r^{2}\Omega^{2}/c^{2}\bigr)^{2}}\right) > 0.
\end{equation}
The key modification with respect to the classical criterion
$d(r^{2}\Omega)^{2}/dr>0$ is that:
\begin{enumerate}
\item The Lorentz factor $\Lf^{2}=(1-v_{\phi}^{2}/c^{2})^{-1}$
  enhances the effective angular momentum at high rotational
  velocities.
\item For a non-barotropic fluid, the full angular momentum
  $\ell$~\eqref{eq:rel-7-9} must be used, and the enthalpy factor
  $(\varepsilon+p)/(\rho_{0}c^{2})$ replaces the trivial factor of
  unity that appears in the non-relativistic case.  A fluid that is
  strongly heated (large $\varepsilon+p$) is thus ``more stable''
  than a cold fluid with the same velocity profile.
\end{enumerate}

\subsection{Relativistic Rayleigh discriminant and epicyclic frequency}
\label{sec:rel-7-66c}

By analogy with Chandrasekhar's non-relativistic
discriminant~(eq.~\eqref{eq:7-18}), we define the
\emph{relativistic Rayleigh discriminant}
\begin{equation}
\label{eq:rel-7-16}
\Phi_{\mathrm{rel}}(r) = \frac{1}{r^{3}\Lf^{4}}\;
  \frac{d}{dr}\!\left(\widetilde{\ell}^{\,2}\right)
= \frac{1}{r^{3}\Lf^{4}}\;
  \frac{d}{dr}\!\left(\Lf^{4}\,r^{4}\Omega^{2}\right).
\end{equation}
Stability requires $\Phi_{\mathrm{rel}}>0$.

The square of the \emph{relativistic epicyclic (radial oscillation)
frequency} is closely related:
\begin{equation}
\label{eq:rel-7-17}
\boxed{\kappa_{\mathrm{rel}}^{2}
= \frac{1}{\Lf^{2}\,r^{3}}\;\frac{d}{dr}\!\left(
  \Lf^{2}\,r^{2}\,v_{\phi}\right)^{2}
= \frac{1}{\Lf^{2}\,r^{3}}\;\frac{d}{dr}\!\left(
  \widetilde{\ell}^{\,2}\right).}
\end{equation}
A displaced fluid ring oscillates radially at frequency
$\kappa_{\mathrm{rel}}$ when the flow is stable
($\kappa_{\mathrm{rel}}^{2}>0$), and grows exponentially when the
flow is unstable ($\kappa_{\mathrm{rel}}^{2}<0$).

In accretion-disk theory (e.g.\ for a test particle in a
Schwarzschild or Kerr metric), $\kappa_{\mathrm{rel}}$ is the
epicyclic frequency that governs the radial structure of density
waves and the location of the innermost stable circular orbit (ISCO).
For a Keplerian disk around a Schwarzschild black hole of mass~$M$,
\begin{equation}
\label{eq:rel-7-18}
\kappa_{\mathrm{Schw}}^{2}
= \Omega_{K}^{2}\!\left(1 - \frac{6GM}{rc^{2}}\right),
\end{equation}
where $\Omega_{K}$ is the Keplerian angular velocity.  The ISCO is at
$r=6GM/c^{2}$ where $\kappa^{2}=0$, precisely corresponding to the
marginal Rayleigh stability of the relativistic circular orbit.

\subsection{Non-relativistic limit}
\label{sec:rel-7-66d}

In the limit $v_{\phi}/c\to 0$, $\Lf\to 1$ and
$\widetilde{\ell}\to r^{2}\Omega$.  Then:
\begin{itemize}
\item The stability condition \eqref{eq:rel-7-14} becomes
  $d(r^{2}\Omega)^{2}/dr>0$, which is Rayleigh's classical
  criterion~\eqref{eq:7-5}.
\item The centrifugal balance \eqref{eq:rel-7-5} reduces to
  $\rho_{0}\,v_{\phi}^{2}/r=-\partial p/\partial r$.
\item The epicyclic frequency \eqref{eq:rel-7-17} becomes
  $\kappa^{2}=(1/r^{3})\,d(r^{2}\Omega)^{2}/dr$, the standard
  result.
\item The relativistic angular momentum \eqref{eq:rel-7-9} reduces
  to $\ell\to r\,v_{\phi}=r^{2}\Omega$.
\end{itemize}
All relativistic modifications vanish, and the entire classical
analysis of Chapter~VII \S\S\,64--66 is recovered.

\subsection{Application to the viscous Couette profile}
\label{sec:rel-7-66e}

For the Couette profile $\Omega(r)=A+B/r^{2}$ (eq.~\eqref{eq:7-1}),
the relativistic stability criterion \eqref{eq:rel-7-14} requires
\begin{equation}
\label{eq:rel-7-19}
\frac{d}{dr}\!\left[\frac{r^{4}(A+B/r^{2})^{2}}
  {\bigl(1-r^{2}(A+B/r^{2})^{2}/c^{2}\bigr)^{2}}\right] > 0
\qquad (R_{1}\le r\le R_{2}).
\end{equation}
When $v_{\phi}\ll c$, the denominator tends to unity and
\eqref{eq:rel-7-19} reduces to $d(Ar^{2}+B)^{2}/dr>0$, yielding
the classical condition $\mu>\eta^{2}$.

At higher velocities the Lorentz-factor weighting makes the criterion
\emph{more restrictive}: the centrifugal ``inertia'' of rapidly
rotating inner fluid is amplified by $\Lf^{4}$, so a steeper rise
in $\Omega$ is needed to maintain stability.  Equivalently,
the critical ratio $\mu_{c}=\Omega_{2}/\Omega_{1}$ at which the
Rayleigh line lies shifts upward:
\begin{equation}
\label{eq:rel-7-20}
\mu_{c}^{\mathrm{rel}} = \eta^{2} + \mathcal{O}
  \!\left(\frac{R_{1}^{2}\Omega_{1}^{2}}{c^{2}}\right).
\end{equation}
The causality constraint $R_{1}\Omega_{1}<c$ ensures that
\eqref{eq:rel-7-19} is always well defined.

\subsection{Causality constraints}
\label{sec:rel-7-66f}

In any physical realisation of relativistic Couette flow, two
causality bounds must be respected:
\begin{enumerate}
\item \textbf{Centrifugal velocity bound:}  $v_{\phi}=r\Omega<c$
  for all $r$ in the flow domain, cf.~\eqref{eq:rel-7-3}.  This is
  an absolute kinematic requirement.
\item \textbf{Signal-speed bound:}  Perturbations propagate at speeds
  bounded by the relativistic sound speed $c_{s}$, with
  $c_{s}^{2}=(\partial p/\partial\varepsilon)_{s}\le c^{2}$.  When
  viscosity is included via the BDNK (Bemfica--Disconzi--Noronha--Kovtun)
  first-order causal theory, the transport coefficients and
  frame coefficients are constrained by coupled inequalities that
  guarantee all viscous and thermal signal speeds remain subluminal
  \citep{Bemfica2018,Kovtun2019}.
\end{enumerate}
These constraints have no analogue in the classical theory and become
essential when assessing stability near the speed-of-light limit.

%% ================================================================
\subsection{Astrophysical application: accretion disk Keplerian flow
  near the ISCO}
\label{sec:rel-7-66-astro}

The relativistic Rayleigh criterion finds its most direct astrophysical
application in the stability of Keplerian accretion disks around compact
objects \citep{NovikovThorne1973,ShakuraSunyaev1973}.  For a test particle
on a circular geodesic in the Schwarzschild spacetime (mass~$M$,
gravitational radius $r_{g}=GM/c^{2}$), the relativistic specific angular
momentum is
\begin{equation}
\label{eq:rel-7-21}
\widetilde{\ell}_{\mathrm{Schw}}^{\,2}(r)
= \frac{r^{2}\,r_{g}}{r-3\,r_{g}}\,,
\end{equation}
which diverges as $r\to 3\,r_{g}$ and attains a minimum at the ISCO,
$r_{\mathrm{ISCO}}=6\,r_{g}$.  At the ISCO,
$d(\widetilde{\ell}^{2})/dr=0$, marking the boundary of Rayleigh
stability: for $r>r_{\mathrm{ISCO}}$ the Keplerian flow is stable
($\Phi_{\mathrm{rel}}>0$), while for
$r<r_{\mathrm{ISCO}}$ no stable circular orbits exist.

For Kerr black holes with spin parameter $a/M$, the ISCO moves inward
as the spin increases: $r_{\mathrm{ISCO}}\approx 6\,r_{g}$ for $a=0$,
$r_{\mathrm{ISCO}}\approx 4.23\,r_{g}$ for $a/M=0.5$, and
$r_{\mathrm{ISCO}}\approx 1.24\,r_{g}$ for $a/M=0.998$
\citep{Bardeen1972}.  In all cases, the Rayleigh stability boundary
coincides with the ISCO, confirming that the relativistic Rayleigh
criterion~\eqref{eq:rel-7-14} correctly predicts the onset of
centrifugal instability in general-relativistic Keplerian flow.

Figure~\ref{fig:rayleigh-isco} illustrates both
$\widetilde{\ell}^{2}(r)$ and its radial derivative for Schwarzschild
and Kerr spacetimes.  The transition from $d\widetilde{\ell}^{2}/dr>0$
(Rayleigh stable) to $d\widetilde{\ell}^{2}/dr=0$ at the ISCO is
clearly visible and shifts to smaller radii with increasing black-hole
spin.

\begin{figure}[htbp]
\centering
\includegraphics[width=0.95\textwidth]{fig_rayleigh_criterion_isco.pdf}
\caption{Left: relativistic specific angular momentum
  $\widetilde{\ell}^{2}$ versus $r/r_{g}$ for Keplerian geodesic flow
  in Schwarzschild ($a/M=0$) and Kerr ($a/M=0.5,\,0.998$) spacetimes.
  Filled markers indicate the ISCO for each case.  Right: radial
  derivative $d(\widetilde{\ell}^{2})/dr$; the Rayleigh-stable region
  ($d\widetilde{\ell}^{2}/dr>0$) is shaded.  The ISCO coincides with
  $d\widetilde{\ell}^{2}/dr=0$ in every case.}
\label{fig:rayleigh-isco}
\end{figure}

\subsubsection{Detailed Kerr analysis: four spin values}

Figure~\ref{fig:isco-rayleigh-kerr} extends the analysis to four
representative spin values $a/M=0$, $0.5$, $0.9$, and $0.998$,
providing a systematic survey of how black-hole spin modifies the
Rayleigh stability of Keplerian flow.  The left panel plots the
specific angular momentum $\ell(r)$ for prograde circular geodesics
in each Kerr spacetime; the minimum of $\ell(r)$ coincides with the
ISCO, whose location shifts dramatically inward with increasing spin:
\begin{equation}
\label{eq:rel-7-22}
r_{\mathrm{ISCO}} \approx
\begin{cases}
6.00\,r_{g} & (a/M=0),\\
4.23\,r_{g} & (a/M=0.5),\\
2.32\,r_{g} & (a/M=0.9),\\
1.24\,r_{g} & (a/M=0.998).
\end{cases}
\end{equation}
The right panel shows the square of the radial epicyclic frequency
$\kappa_{r}^{2}(r)$, which vanishes at the ISCO and becomes negative
(Rayleigh-unstable) inside it.  This confirms the identification of
the ISCO with the onset of centrifugal instability for all spins.

The behaviour near extreme spin ($a/M\to 1$) is of particular interest
for EHT observations \citep{EHT2019M87V} and GRMHD simulations
\citep{Porth2019}.  For $a/M=0.998$, the ISCO lies at
$r\approx 1.24\,r_{g}$, barely outside the event horizon at
$r_{+}\approx 1.06\,r_{g}$.  In this regime the radial epicyclic
frequency $\kappa_{r}$ achieves a pronounced maximum before dropping
to zero, creating a narrow radial zone where trapped oscillations
(precursors to QPOs) can form.

Recent GRMHD simulations of magnetically arrested disks
\citep{Ripperda2022} show that magnetic flux eruptions near the ISCO
are closely tied to the vanishing of $\kappa_{r}^{2}$: the loss of
radial centrifugal support allows accumulated magnetic flux to be
violently ejected, producing observable flares.  The Rayleigh
stability analysis of this section provides the hydrodynamic
foundation for understanding these phenomena.

\begin{figure}[htbp]
\centering
\includegraphics[width=0.95\textwidth]{fig_isco_rayleigh_kerr.pdf}
\caption{Rayleigh criterion near the ISCO for Kerr black holes.
  Left: specific angular momentum $\ell(r)$ for prograde geodesics
  with $a/M=0$ (Schwarzschild), $0.5$, $0.9$, and $0.998$.  Filled
  circles mark the ISCO where $d\ell/dr=0$.
  Right: radial epicyclic frequency $\kappa_{r}^{2}(r)$ for the same
  four spins.  The Rayleigh-unstable region ($\kappa_{r}^{2}<0$,
  shaded) lies inside the ISCO.  As spin increases, the ISCO moves
  inward and $\kappa_{r}$ develops a sharper maximum, creating
  a narrow zone for trapped oscillations.}
\label{fig:isco-rayleigh-kerr}
\end{figure}

%% ================================================================
\section{Exact stationary Couette-type solutions: the Gavassino--Niekamp--Schlichting--Denicol theory}
\label{sec:rel-7-gavassino-couette}
%% ================================================================

The inviscid analysis of the preceding sections establishes the Rayleigh
criterion for relativistic circular flows.  However, the base state around
which one perturbs is itself modified once dissipation is included.
\citet{GavassinoCouette2025} (Gavassino, Niekamp, Schlichting \& Denicol,
arXiv:2512.10420) have recently derived \emph{exact} stationary solutions for
relativistic dissipative Couette-type flows---the relativistic generalization
of the classical planar Couette profile.  Their results reveal that the
``inertia of heat'' (the momentum carried by the heat flux) qualitatively
alters the velocity profile, a purely relativistic effect with no Newtonian
counterpart.

\subsection{Setup and the role of heat conduction}
\label{sec:rel-7-gavassino-setup}

Consider a viscous relativistic fluid confined between two infinite parallel
plates at $x = \pm L/2$, moving in opposite directions with velocities
$v_y = \pm v$, and maintained at temperatures $T_{\pm}$.  In
the classical (non-relativistic) theory, the velocity profile is simply
linear: $u^y = 2vx/L$.  \citet{Rogava1996} derived a relativistic
generalization by setting the heat conductivity to zero, obtaining
\begin{equation}
\label{eq:rel-7-rogava}
u^y_{\text{Rogava}} = \frac{2vx}{L\sqrt{1-v^2}}\,.
\end{equation}
However, Gavassino et al.\ show that setting the heat conductivity $\chi$ to
zero is inconsistent: viscous heating necessarily warms the fluid interior,
and a steady state requires finite thermal conductivity to evacuate the
excess energy through the boundaries.

\subsection{The exact velocity profile}
\label{sec:rel-7-gavassino-velocity}

Working in the Eckart frame (where the four-velocity is aligned with the
particle current), Gavassino et al.\ show that the Israel--Stewart
relaxation terms vanish identically for stationary planar-symmetric
flows (since $u^\mu \partial_\mu = 0$ when $u^x = 0$).  The momentum
conservation equation $\partial_x T^{xy} = 0$ gives
\begin{equation}
\label{eq:rel-7-gavassino-ode}
\eta \,\partial_x u = \gamma u \bar{A} - \gamma^2 \bar{C}\,,
\end{equation}
where $\bar{A} = T^{xt}/\eta$ and $\bar{C} = T^{xy}/\eta$ are
constants determined by boundary conditions.
The crucial point is the term $\gamma u \bar{A}$, which arises from the
\emph{inertia of heat}: the heat flux $q^x$ carries momentum, contributing
to $T^{xy}$ through the coupling $u q^x$.  This is absent in
Rogava's approximation (which sets $q^x = 0$).

For equal boundary temperatures ($T_+ = T_-$), symmetry forces $\bar{A}=0$,
and the ODE reduces to $\partial_x u = -(1+u^2)\bar{C}$, which integrates
to the \emph{exact} velocity profile:
\begin{equation}
\label{eq:rel-7-gavassino-exact}
\boxed{u(x) = \tan\!\left[\frac{2x}{L}\,
    \arctan\!\left(\frac{v}{\sqrt{1-v^2}}\right)\right].}
\end{equation}
This is a remarkable result: despite the complexity of relativistic
dissipative hydrodynamics, the velocity profile is given by a single
transcendental function.  Key properties:
\begin{enumerate}
\item In the non-relativistic limit $v \ll 1$, $\arctan(v/\sqrt{1-v^2})
  \to v$, and we recover the classical linear profile $u \approx 2vx/L$.
\item In the ultrarelativistic limit $v \to 1$, $u(x) \to \tan(\pi x/L)$,
  which remains finite throughout the domain (diverging only at the
  boundaries), whereas Rogava's approximation \eqref{eq:rel-7-rogava}
  diverges everywhere except at $x=0$.
\item The profile is \emph{independent of the thermal conductivity} $\chi$
  (though it requires $\chi \neq 0$ for consistency).  This universality
  holds whenever $\eta$ is temperature-independent.
\end{enumerate}

\subsection{Temperature and heat flux profiles}

The associated temperature and heat flux are also obtained analytically:
\begin{equation}
\label{eq:rel-7-gavassino-temp}
T(x) = T_+ + \frac{\eta}{\chi}\,\ln\!\left[
  \frac{1}{\sqrt{1-v^2}}\,\cos\!\left(
    \frac{2x}{L}\,\arctan\!\left(\frac{v}{\sqrt{1-v^2}}\right)
  \right)\right],
\end{equation}
\begin{equation}
\label{eq:rel-7-gavassino-heat}
q^x(x) = \frac{2\eta}{L}\,
  \arctan\!\left(\frac{v}{\sqrt{1-v^2}}\right)\,
  \tan\!\left[\frac{2x}{L}\,
    \arctan\!\left(\frac{v}{\sqrt{1-v^2}}\right)\right].
\end{equation}
The temperature peaks at the centre (maximum viscous heating) and the
heat flux is directed outward toward the plates, as expected physically.

\subsection{Implications for perturbation analysis}
\label{sec:rel-7-gavassino-implications}

The Gavassino et al.\ exact solution has profound implications for the
stability analysis of relativistic Couette flow developed in the earlier
sections of this chapter:

\begin{enumerate}
\item \textbf{Modified base state.}  The inviscid analysis of
  \S\ref{sec:rel-7-67} uses Rogava's profile \eqref{eq:rel-7-rogava} as
  the base state.  Gavassino et al.\ show that the correct dissipative
  base state is \eqref{eq:rel-7-gavassino-exact}, which has a qualitatively
  different curvature: the $\tan$ profile develops larger gradients near
  the boundaries than the linear (or Rogava) profile, potentially
  concentrating perturbation energy near the plates.

\item \textbf{Inertia of heat in the Landau frame.}  In the Landau frame
  (where $u^\mu$ tracks energy flow), the heat flux causes the velocity
  to acquire a component $u_L^x \neq 0$ perpendicular to the plates, with
  \begin{equation}
  \label{eq:rel-7-landau-deflection}
  \frac{u_L^x}{u_L^y} = \frac{2\eta}{L(\varepsilon+P)}\,
    \arctan\!\left(\frac{v}{\sqrt{1-v^2}}\right)
    \xrightarrow{v\to 1} \frac{\pi\eta}{L(\varepsilon+P)}
    \sim \text{Kn}\,.
  \end{equation}
  This transverse deflection of order the Knudsen number is a genuine
  relativistic frame effect that must be accounted for when computing
  perturbation growth rates in the Landau frame.

\item \textbf{Asymmetric profiles from temperature gradients.}
  When $T_+ \neq T_-$, the heat-inertia coupling breaks the odd
  symmetry of the velocity profile.  The general solution involves
  trigonometric and hyperbolic functions (Eq.~(45) of
  \citealt{GavassinoCouette2025}), with the asymmetry growing with
  $v$ since the inertia-of-heat term $u q^x$ increases with velocity.

\item \textbf{Robustness of the Israel--Stewart regulator.}
  A key technical finding is that the IS relaxation terms vanish
  identically for stationary planar flows (because $u^\mu\partial_\mu = 0$
  in the Eckart frame).  The dissipative solution is therefore independent
  of the relaxation times $\tau_\Pi, \tau_q, \tau_\pi$, making it
  a robust prediction shared by all causal dissipative frameworks
  (IS, BDNK, DNMR).
\end{enumerate}

\medskip\noindent\textbf{Connection to regime of applicability.}
Wagner \& Gavassino (arXiv:2309.14828, \citealt{WagnerGavassinoRegime2024})
have systematically assessed when Israel--Stewart hydrodynamics is
trustworthy versus when first-order (BDNK/Navier--Stokes) theories
suffice.  Their ranking of IS truncations in decreasing accuracy---IReD,
DNMR, second-order gradient expansion, 14-moment---provides a criterion
for each section of this book: the IS formulation is falsifiable, the
relaxation time is observable, and the IReD prescription for transport
coefficients ($\tau_\Pi = \sum_n \zeta_n \tau_n / \sum_m \zeta_m$) is
superior to DNMR whenever the slowest non-hydrodynamic mode has small
coupling to the relevant viscous flux.
